\documentclass[10pt]{article}
\usepackage[margin=.75in]{geometry}
\usepackage{amsmath,amssymb,amsthm,mathtools,booktabs,array,hyperref}
\hypersetup{hidelinks}
\newtheorem{theorem}{Theorem}
\newtheorem{lemma}{Lemma}
\newcommand{\half}{\tfrac12}
\newcommand{\bc}{b_{\mathrm c}}
\newcommand{\Zc}{Z_*}
\newcommand{\Qh}{Q_*}
\newcommand{\Rh}{R_*}
\newcommand{\Bh}{B_*}
\newcommand{\D}{\mathcal D}
\newcommand{\Ps}{\Psi}
\newcommand{\E}{\mathcal E}
\newcommand{\W}{\mathcal W}
\newcommand{\Aop}{\mathcal A}
\newcommand{\Bop}{\mathcal B}
\newcommand{\calP}{\mathcal P}
\newcommand{\ps}{\succeq}
\setlength{\jot}{3pt}
\title{A uniform contact certificate with independent source and testing parameters}
\author{}\date{}
\begin{document}
\maketitle\vspace{-2em}

The Bell slack is quadratic in the state coordinates. At a zero, every
square in a positive decomposition must vanish. To retain this zero
throughout the source family, we combine two square completions before
estimating their singular source coefficients.

Write the source as $c=(d,\lambda,n,m,C,A,D,\alpha,\beta)\geq0$ and
put $D_u=1+u^2$, $D_v=1+v^2$. The canonical Bell slack $S_c=F_h/4$ is
\begin{align}
 D_uD_vS_c={}&\lambda D_u(W+vX)^2+nD_uD_vX^2
 +m(vW-X+uvY-uZ)^2 \notag\\
 &+CD_v(uW-Y)^2+AD_uD_vY^2
 +(D+d/2)(uW+uvX-Y-vZ)^2 \notag\\
 &+(\alpha+d/2)D_u(vY-Z)^2
 +(\beta+d/2)D_v(X+uZ)^2-\tfrac d2D_uD_vZ^2 .
 \label{eq:slack}
\end{align}
All four state coordinates are independent. Use the testing change
\begin{equation}
 t=\frac{X+uZ}{D_u},\quad r=\frac{Y-uW}{D_u},\quad
 w=\frac{W+uY}{D_u},\quad \zeta=\frac{Z-uX}{D_u},
 \qquad
 (W,X,Y,Z)=(w-ur,t-u\zeta,r+uw,\zeta+ut).
 \label{eq:rotation}
\end{equation}
Only the strictly positive testing factor $D_u$ is inverted.

The source family in this note has $(d,\lambda,\beta)=(1,0,0)$.
Its rational coefficients are given in Appendix~\ref{app:source} in terms of
\begin{equation}
 \delta=2^{-10},\qquad
 0<B,T\leq\delta,\quad0<x=jB\leq\delta,\quad
 0<p<\min(j,1),\quad0<H<1,\quad n\geq0.
 \label{eq:domain}
\end{equation}
The condition $n\geq0$ is imposed on the displayed rational coefficient.
There is no upper or lower bound on $j$ or $jT$.

\begin{theorem}\label{thm:cap}
Every source in \eqref{eq:domain} satisfies $S_c\geq0$ for every state and
\[
 |u|\leq u_0=2^{-12},\qquad |v|\leq v_0=2^{-25}.
\]
On each required real sign chart the construction uses finitely many fixed
testing functions and source-only positive Gram matrices. It includes
$n=0$ literally and allows arbitrary nonnegative additions of the eight
event coordinates. Source and testing parameters remain independent.
\end{theorem}

The common source weight
\[
 \phi=\frac{p^2}{1-p+p^2}
\]
makes the two available completions compatible. The factor $1-\phi$ cancels
$1-p$ in one denominator; $\phi$ controls the other denominator through a
harmonic coefficient. The first-normal expression $J_\alpha$ below is kept
whole while this allocation is made.

\section{One weight for the two denominators}
Put $y=TB$, $q_z=jT$, $\mu=m(1-p)$ and define
\begin{gather}
 N=t-pvw,\quad M=-r-Hv\zeta,\quad
 \bc=vw-u\zeta/j,\quad R=\zeta-uv\,w/B,\quad R_u=v(u^2-x)w,\notag\\
 h_0=\mu+\half,\quad
 M_e=\frac{nh_0j^2}{n+h_0},\quad
 \widehat N=\sqrt{n+h_0}\left(N+\frac{nj}{n+h_0}\bc\right),\quad
 E_*=\sqrt{M_e}\bc .
 \label{eq:normal}
\end{gather}
Completing the square gives
$h_0N^2+n(N+j\bc)^2=\widehat N^2+E_*^2$.
The positive source masses $f_c,q_c,r_c$ in Appendix~\ref{app:source} satisfy
\begin{equation}
 \Sigma=M_e+f_c/y>0,\qquad p^2\leq32\Sigma,\qquad
 \theta=M_e/\Sigma\in[0,1].
 \label{eq:harmonic}
\end{equation}
Appendix~\ref{app:source} proves this estimate without a bound on $j$.
Consequently, with $D_p=1-p+p^2\geq3/4$,
\begin{equation}
 \frac{1-\phi}{1-p}=\frac1{D_p}\leq\frac43,\qquad
 \frac{\phi}{\Sigma}\leq\frac{128}{3},\qquad
 \frac{\phi}{j}\leq\frac{4p}{3}\leq\frac43 .
 \label{eq:phi}
\end{equation}
Every quantity in \eqref{eq:phi} depends only on the source.

To identify the corresponding squares, put
\begin{gather}
 U=u/\sqrt x,\quad V=v/\sqrt y,\quad \Zc=\zeta/\sqrt{q_z},\quad
 \Qh=Uw-V\Zc,\quad \Bh=Vw-U\Zc,\quad\Rh=\Zc-UVw,\notag\\
 \Ps=(U-V)^2w^2+\Zc^2(1-UV)^2+2UV(w-\Zc)^2 .
 \label{eq:circuit}
\end{gather}
Here $\bc=\sqrt y\,\Bh$ and $R=\sqrt{q_z}\Rh$.
The first-normal term and its exact decomposition are
\begin{equation}
 J_\alpha=j\alpha\bc=-\frac{\alpha}{B}R_u-\alpha uR .
 \label{eq:J}
\end{equation}
Let $A_1=(\alpha/B)R_u$, $B_1=\alpha uR$. Since
$(1+U^2)\Ps-\Bh^2=\Rh^2+U^2\Qh^2+U^2(w-UV\Zc)^2$,
\begin{align}
 &2(1-\phi)(\alpha/B)^2R_u^2
 +2(1-\phi)\alpha^2q_zu^2\Rh^2
 +\phi j^2\alpha^2y(1+u^2/x)\Ps-J_\alpha^2 \notag\\
 &\qquad=(1-\phi)(A_1-B_1)^2
 +\phi j^2\alpha^2y[\Rh^2+U^2\Qh^2+U^2(w-UV\Zc)^2].
 \label{eq:Jpurchase}
\end{align}
For $\eta=B+x$, the three coefficients on the left, relative respectively to
$R_u^2$, $r_c\Rh^2$ and $f_c\Ps$, are at most
\[
 2\eta^2,\qquad (64/3)\eta^2u^2,\qquad
 8B\eta^2+(32/3)\eta^2u^2 .
\]
Thus keeping $J_\alpha$ whole adds only a small coefficient to
$R_u^2$, even as $p\to1$. The same $\phi$ gives bounds for $R_u$ and the second-normal
term
\begin{equation}
 O=\sqrt T\,v(\alpha-BTCHu^2)R,\qquad
 H_1=pv^2\bc,\qquad H_2=\sqrt T\,puv\bc .
 \label{eq:helpers}
\end{equation}
The exact identities and their simultaneous bounds are given in
Appendix~\ref{app:purchases}; none introduces a moving testing radical.

\section{Both normals in one positive matrix}
Set
\[
 a=(n+h_0)^{-1/2},\quad
 b_0=\sqrt{\frac{n}{h_0(n+h_0)}},\quad
 \tau=\sqrt{TH(1-H)},\quad \rho_Q=T\sqrt{x/q_c},
\]
\[
 \widehat M=M/\tau,\quad
 \widehat Q=(v\zeta-Tuw)/\rho_Q=-\sqrt{q_c}\Qh,\quad
 \widehat R=\sqrt{r_c}\Rh,\qquad N=a\widehat N-b_0E_*.
\]
In particular $a^2+b_0^2=1/h_0\leq2$. Choose the three explicit factors
$\psi_i$ of $f_c\Ps$ from \eqref{eq:circuit} and use the thirteen rows
\[
 f=(\widehat N,E_*,R_u,J_\alpha,H_1,\widehat M,O,H_2,
                         \widehat Q,\widehat R,\psi_1,\psi_2,\psi_3)^T.
\]
The quadratic form $f^TKf$ is specified by its diagonal
\begin{equation}
 (1,\half,16,16,8,a_M,16,8,\half,\half,\tfrac14,\tfrac14,\tfrac14),
 \qquad a_M=\tau^2(C+D+\half+A),
 \label{eq:diag}
\end{equation}
and the following complete list of cross terms:
\begin{align}
 &2N(K_RR_u-J_\alpha+H_1/2)
 +2\sqrt{H(1-H)}\,\widehat M\{O+(\mu/p)H_2\}\notag\\
 &\quad+2uvN\{(K_c/j)\rho_Q\widehat Q
                         +(\alpha+\half)\tau\widehat M\}
 +2\tau\rho_Q C_q\,\widehat M\widehat Q, \label{eq:cross}\\
 K_c&=p/2+j[H(\alpha+\half)-\half]+n(p-j),\notag\\
 K_R&=-\mu-(\alpha+\half)(1-p)+(T/j)K_c,\qquad
 C_q=A/T+Hv^2[C(1+u^2)+A+\alpha+\half].\notag
\end{align}

The source bounds give $|K_R|<1$, $\mu/p\leq2/3$ and $1/8<a_M<3/4$.
For the constant matrix excluding the constant $\widehat M\widehat Q$
entry, subtracting $I/16$ gives two Schur corrections at most
\[
 2(32/255+4/127)=10168/32385<7/16,\qquad
 4/255+16/1143<1/16 .
\]
Its remaining rows are diagonal. It therefore dominates $I/16$.
The excluded entry has norm at most $T$, and the exact support is
\begin{equation}
 K=K_{00}+uvK_{11}+v^2K_{02}+u^2v^2K_{22},\qquad
 \|K_{11}\|,\|K_{02}\|,\|K_{22}\|\leq1.
 \label{eq:Kbounds}
\end{equation}
The derivation is in Appendix~\ref{app:source}.
Let $\Pi_{11}$ be the coordinate projection onto
$(\widehat N,E_*,\widehat M,\widehat Q)$, and let
$\Pi_{02}=\Pi_{22}$ project onto $(\widehat M,\widehat Q)$.
These projections contain the exact supports of the three nonconstant
coefficient matrices. The fixed monomials $P_{ab}=(u/u_0)^a(v/v_0)^b$, for
$(a,b)=(1,1),(0,2),(2,2)$, have budgets
$2^{-37},2^{-50},2^{-74}$. The identity
\[
 K=K_{00}-\sum\rho_{ab}\Pi_{ab}+
 \sum_{a,b}\left[
 \frac{\rho_{ab}\Pi_{ab}+A_{ab}}2(1+P_{ab})+
 \frac{\rho_{ab}\Pi_{ab}-A_{ab}}2(1-P_{ab})\right],
 \quad A_{ab}=u_0^av_0^bK_{ab},\quad\rho_{ab}=u_0^av_0^b
\]
has positive source matrices and a constant reserve greater than $I/32$,
because $1/16-\delta-2^{-37}-2^{-50}-2^{-74}>1/32$.

\section{Restoring the scalar squares}
The complete parent identity, written out in Appendix~\ref{app:parent}, is
\begin{equation}
 D_vS_c=f^TKf+\calP+
 \tfrac34 f_c\Ps+\half q_c\Qh^2+\half M_e\bc^2+\half r_c\Rh^2
       -g_c\Ps-2\kappa V^2\Qh^2-2\kappa\Rh^2-\mathcal T ,
 \label{eq:coreidentity}
\end{equation}
where $\kappa=-k_c>0$, $\calP$ is a displayed sum of retained positive
squares, and
$\mathcal T=16R_u^2+16J_\alpha^2+16O^2+8H_1^2+8H_2^2$.
The identities in Appendix~\ref{app:purchases} replace all negative terms
in \eqref{eq:coreidentity} by positive remainders. Their total fractions
$b_i$ of $(f_c\Ps,q_c\Qh^2,M_e\bc^2,r_c\Rh^2)$ satisfy
\begin{equation}
 \|b_\Ps\|_*<9/16,\qquad \|b_Q\|_*<2^{-22},\qquad
 \|b_E\|_*<1/16,\qquad \|b_R\|_*<1/64 .
 \label{eq:totalbounds}
\end{equation}
The norm sums absolute coefficients weighted by $u_0^av_0^b$.
The core uses fractions $(1/4,1/2,1/2,1/2)$, so each remaining scalar
coefficient is strictly positive with a fixed coefficient reserve.

More explicitly, write $b_i=c_{i0}+\sum_\nu c_{i\nu}P_\nu$, where the
nonconstant exponents lie in
\[
 (2,0),(4,0),(0,2),(2,2),(4,2),(6,2),(0,4),(2,4),(4,4).
\]
Since these exponents are even, $0\leq P_\nu\leq1$ on the signed cap and
$P_\nu$ is a polynomial square. With $c^\pm=\max(\pm c,0)$,
\[
 1-d_i-b_i=
 \underbrace{1-d_i-c_{i0}-\sum_\nu c_{i\nu}^+}_{>0}
 +\sum_\nu c_{i\nu}^+(1-P_\nu)+\sum_\nu c_{i\nu}^-P_\nu .
\]
This uses only fixed scalar gates. Expanding the finitely many displayed
roots separates their source coefficients from their testing monomials.
Appendix~\ref{app:functions} records the resulting fixed lists.

All expressions are literal at $n=0$: $b_0=M_e=E_*=0$ and
$\Sigma=f_c/y>0$. On $u,v\geq0,w\zeta\leq0$, replace the three factors
of $f_c\Ps$ by the five factors
$\sqrt{f_c}(Uw,Vw,\Zc,UV\Zc,2\sqrt{-UVw\Zc})$.
Only the three uncoupled diagonal rows change to five rows with the same
coefficient $1/4$; every matrix bound remains valid.
Appendix~\ref{app:parent} gives positive expansions of both radial forms
on this state chart. The source-preserving angle gauges extend the two
identities to all angle signs. All testing roots have their literal zero
values on the axes; the independent core completion \eqref{eq:signcore}
also gives direct nonnegativity there.
Adding nonnegative event coordinates in \eqref{eq:slack} preserves positivity.
Complex states are covered by their independent real and imaginary parts.

The theorem concerns this source corner and this fixed testing cap.
Complementary testing regions and admission of all canonical sources require
further arguments.

\appendix
\section{Source coefficients and their uniform bounds}\label{app:source}
The source is the stationary RSR dictionary of the proof companion,
RSR4--RSR6 and RSR15--RSR19. We give its polynomials here so that the
bounds use only \eqref{eq:domain}. The symbol $F$ in this appendix is a
source polynomial, distinct from the homogeneous slack $F_h$.
Set $x=jB$, $y=TB$, $q_z=jT$, and define
\input{source_dictionary_exact}
The condition $n\geq0$ means $\mathsf N\geq0$.
Every denominator in this dictionary is positive on \eqref{eq:domain}.
In particular $1-HT>0$,
$1-y[(1-H)j-p]\geq1-Tx>0$ and $E\geq1-\delta>0$.

Put $J=j+1$, $K=J-p$, and
\[
 f_0=1-\delta-\delta^3,\quad
 \ell_0=1+\delta+\delta^2+\delta^3,\quad
 R_0=1+2\delta+2\delta^2+4\delta^3+\delta^4(1+2\delta).
\]
Because $p<\min(j,1)$, $J/2\leq K<J$ and $K>1$.
Termwise bounds in the displayed polynomials give
\begin{equation}
 J\leq P\leq R_0J,\qquad
 1-\delta\leq E\leq1+\delta^2,\qquad
 f_0K\leq L\leq\ell_0K,\qquad f_0K\leq F\leq2K.
 \label{eq:fullsourcebounds}
\end{equation}
Here $P\geq J$ follows from its positive terms; in the upper bound use
$BJ=B+x\leq2\delta$ to bound its successive corrections, divided by $J$,
by $2\delta,2\delta^2,4\delta^3,\delta^4(1+2\delta)$.
For $L$, use $W_0=K-Hj$ to obtain
\[
 L\geq K-T-BT^2\geq f_0K,\qquad
 L\leq K[1+T+BT+BT^2]\leq\ell_0K.
\]
For $F$, the exact regrouping
\begin{align*}
 F_1={}&2Bjp(1-H)+BHj+2Bp(1-p)+(1-H)J-Hp+p(1-p)/j,\\
 F_2={}&B^2Hjp(1-H)+B^2Hp(1-p)+BHj(1-H)-BH^2p+BHp(1-p)/j
\end{align*}
gives $F\geq K-Tp-BT^2p\geq f_0K$.
For the upper bound,
$F_1\leq(3+5\delta)K$ and $F_2\leq(2\delta+2\delta^2)K$,
which yield $F<2K$ at $\delta=2^{-10}$.

Write
\[
 G_q=\frac{(1+x)(1+Hy)(K+Hjy)}{1-HT}-xF,\qquad
 q_c=\frac{TxG_q}{2P}.
\]
The first term in $G_q$ is at least $K$, so $G_q\geq(1-2\delta)K$.
The dictionary consequently gives
\begin{equation}
 f_c/y>j^2B/5,\quad
 f_c\geq Tx^2/8,\quad q_c\geq Tx/8,\quad
 r_c\geq TjB^2(1-p)/8,\quad
 \alpha\leq B(B+x),\quad c_0:=BTCH<B .
 \label{eq:positivesourcebounds}
\end{equation}
For the first four inequalities use, respectively,
\[
 \frac{f_0}{4R_0}>\frac15,\qquad
 \frac{f_0}{4R_0}>\frac18,\qquad
 \frac{1-2\delta}{4R_0}>\frac18,\qquad
 \frac{f_0}{4(1+\delta^2)R_0}>\frac18 .
\]
For $\alpha$, bound $(K+Hjy)/P\leq1+\delta^2$ and $L\leq\ell_0J$:
\[
 \alpha\leq B^2J\,\frac{(1+\delta^2)\ell_0}{2(1-\delta)}
 \leq B(B+x).
\]
The last scalar factor is less than one. Also $c_0=BL/(2P)<B$.
Directly,
\begin{equation}
 \mu/p\leq(1+\delta)(1+\delta^2)^2/2<2/3,\qquad
 A/T\leq\frac{(1+\delta)(1+\delta^2)^2}{2(1-\delta)}<2/3 .
 \label{eq:mubounds}
\end{equation}

\subsection*{The harmonic bound}
Put $\rho=p/j<1$. Bounding the negative terms of $\mathsf N/j$ below
and discarding only nonnegative terms gives
\begin{gather}
 \mathsf N/j\geq\rho-BD_N,\qquad
 D_N=1+2j+T(1+j)+BT(2j+4j^2)+BT^2j(1+Bj),\notag\\
 D_N/J\leq2+\delta+4\delta^2+\delta^3(1+\delta)<9/4.
 \label{eq:Nlower}
\end{gather}
Since $\mathsf N/j=2(1-\rho)Pn$ and $\rho J=p+\rho<2$,
\[
 \rho^2\leq\frac{2\rho}{J}
 \leq4(P/J)n+2B(D_N/J)
 \leq5n+\tfrac92B .
\]
If $n\leq h_0$, $M_e\geq nj^2/2$, and hence
\[
 p^2\leq5nj^2+\tfrac92Bj^2
 \leq10M_e+\tfrac{45}{2}f_c/y
 \leq\tfrac{45}{2}\Sigma .
\]
If $n\geq h_0$, then $M_e\geq h_0j^2/2\geq j^2/4$ and $p^2<j^2\leq4M_e$.
This proves \eqref{eq:harmonic}. At $n=0$ the positive term $f_c/y$
is unchanged.

\subsection*{Signed coefficients and the normal matrix}
The exact positive expression
\begin{align*}
 \mathsf K-pW_0+Bp(1-p)J
 ={}&B(1-p)[Hj(p+1)+H(1-p)+p^2]\\
 &+TBH[(p+1)(1-H)j+(1-p)(BHj+p+1-H)]\\
 &+T^2B^2H^2j(1-H)
\end{align*}
shows $\mathsf K>0$, since $BJ<1$ and $W_0\geq1-p$.
For an upper bound, its four displayed lines contribute at most
$p,B,B\delta(2+\delta),B\delta^3$ after division by $J$.
Thus $\mathsf K/J\leq p+3B$.
Thus $\kappa=-k_c>0$; the exact ratios imply
\begin{equation}
 \phi\kappa/f_c\leq3T,\quad
 \phi\kappa/(\Sigma yx)\leq24T,\quad
 \kappa/(q_cy)\leq3,\quad
 (1-\phi)\kappa/r_c\leq3T,\quad
 g_c/f_c=HyL/F\leq2y.
 \label{eq:kappabounds}
\end{equation}
For example the four rational constants proving the first four bounds are
bounded, respectively, by
\[
 \frac{8(1+\delta)(1+3\delta)}{3(1-\delta)f_0}<3,\quad
 \frac{64(1+\delta)(1+3\delta)}{3(1-\delta)}<24,\quad
 \frac{(1+\delta)(2+6\delta)}{(1-\delta)(1-2\delta)}<3,\quad
 \frac{4(2+6\delta)}{3f_0}<3.
\]
They use \eqref{eq:phi}, $K\geq J/2$ and the exact source ratios.

For the first normal, bounding the positive terms in $\mathsf N/j$ above
and discarding its nonpositive terms gives
\[
 \mathsf N/j\leq1+4\delta+2\delta^2+6\delta^3+
                       \delta^4(1+\delta)<2.
\]
In this bound use $Hj+H-j-1\leq0$ and
$2Hjp-2jp\leq0$ before bounding the other terms.
As $P\geq J$, it follows that
$0\leq n(1-p/j)=\mathsf N/(2jP)<1$.
Therefore $|K_c/j|\leq2$, and
\[
 |K_R|\leq \mu+(\alpha+\half)(1-p)+2T
 \leq\tfrac23+2\delta^2+2\delta<1.
\]
For the second normal, the complementary $H$ and $1-H$ factors in $C,D$
must be retained together:
\[
 a_M=TH(1-H)(C+D+\half+A),\qquad
 \frac{f_0}{4R_0}\leq a_M
 \leq\frac{\max\{\ell_0,(1+\delta^2)^3\}}2+\frac{\delta}{8}
                              +\frac{\delta^2}{6}.
\]
The endpoints lie strictly between $1/8$ and $3/4$.
Indeed $THC=L/(2P)$, while
$T(1-H)D=(1+Hy)(K+Hjy)[1-y((1-H)j-p)]/(2P)$.
Their lower bounds are $f_0/(4R_0)$ and
$(1-\delta^2)/(4R_0)$, and their upper bounds are
$\ell_0/2$ and $(1+\delta^2)^3/2$. The terms $\tau^2/2$ and $\tau^2 A$
contribute at most $\delta/8$ and $\delta^2/6$.

Finally $\rho_Q\leq\sqrt{8T}$ and $\tau\leq\sqrt T/2$.
Reading the three cross coefficients in \eqref{eq:cross} gives
\begin{gather*}
 \|K_{11}\|^2\leq578T/9<1,\\
 \|K_{02}\|\leq\sqrt2[\ell_0/2+
                  \delta(2\delta/3+2\delta^2+\half)]<1,\qquad
 \|K_{22}\|\leq\sqrt2\,\ell_0/2<1 .
\end{gather*}
The constant $\widehat M\widehat Q$ entry is at most
$\tau\rho_QA/T\leq2\sqrt2\,T/3<T$.
This proves every matrix bound used in the main argument.


\section{Exact identities and simultaneous coefficient bounds}\label{app:purchases}
In this appendix $s=u^2$ and $\eta=B+x$. We compare positive multiples
of the four forms
\[
 (f_c\Ps,\ q_c\Qh^2,\ M_e\bc^2,\ r_c\Rh^2).
\]
The fractions $(b_\Ps,b_Q,b_E,b_R)$ bound a row when
the sum of these four weighted forms minus its square is an explicitly
displayed sum of squares.

The elementary identities needed below are
\begin{align}
 \Aop:=(1+U^2)\Ps-\Rh^2
   &=\Bh^2+U^2\Qh^2+U^2(w-UV\Zc)^2,\notag\\
 \E:=(1+V^2)\Ps-V^2\Bh^2
   &=\Qh^2+\Rh^2+V^2(w-UV\Zc)^2 . \label{eq:elementary}
\end{align}
There is also the opposite form
\begin{align}
 \Bop:=2(\Ps+V^2\Qh^2+U^2\Bh^2)-\Rh^2
   &=(E_0+T_0)^2+2(U-V)^2w^2+4UV(w-\Zc)^2\notag\\
   &\quad+2J_0^2+4UV(U-V)^2\Zc^2,\notag\\
 E_0&=\Zc(1-UV),\quad T_0=UV(w-\Zc),\quad
 J_0=T_0-(U-V)^2\Zc. \label{eq:opposite}
\end{align}
All are sums of squares for $U,V\geq0$.

\subsection*{A common bound for the opposite row}
For a nonnegative polynomial multiplier $a_0$, define the fractions
\begin{align}
 \mathfrak b_\Ps(a_0)&=a_0\phi\left[\frac{1+\theta}{f_c}
                                      +\frac{s}{\Sigma yx}\right],&
 \mathfrak b_Q(a_0)&=\frac{2a_0\phi\theta v^2}{q_cy},\notag\\
 \mathfrak b_E(a_0)&=\frac{2a_0\phi s}{\Sigma yx},&
 \mathfrak b_R(a_0)&=\frac{a_0(1-\phi)}{r_c}.
 \label{eq:oppositefractions}
\end{align}
Since $f_c/(\Sigma y)=1-\theta$ and $M_e/\Sigma=\theta$, exactly
\begin{align}
 &f_c\mathfrak b_\Ps(a_0)\Ps+q_c\mathfrak b_Q(a_0)\Qh^2
 +M_e\mathfrak b_E(a_0)\bc^2+r_c\mathfrak b_R(a_0)\Rh^2-a_0\Rh^2
 \notag\\
 &\hspace{30mm}=a_0\phi[(1-\theta)\Aop+\theta\Bop].
 \label{eq:oppositepurchase}
\end{align}
The direct fraction on $\Rh^2$ is weighted by $1-\phi$ before
$r_c^{-1}$ is estimated.

For $O=\sqrt{j}\,Tv(\alpha-c_0s)\Rh$, put
\[
 \bar a=2jT^2v^2(\alpha^2+c_0^2s^2),\qquad c_0=BTCH.
\]
There is the exact additional square
\[
 \bar a\Rh^2-O^2=jT^2v^2(\alpha+c_0s)^2\Rh^2 .
\]
Therefore \eqref{eq:oppositepurchase} with $a_0=\bar a$ bounds $O^2$.
Using \eqref{eq:phi} and \eqref{eq:positivesourcebounds},
its coefficientwise absolute majorants are
\begin{equation}
 \begin{split}
 \mathfrak b_\Ps(\bar a)&\leq_*
             (128/3)Tv^2(1+2s)(\eta^2+s^2),\\
 \mathfrak b_Q(\bar a)&\leq_*32v^4(\eta^2+s^2),\\
 \mathfrak b_E(\bar a)&\leq_*(512/3)Tsv^2(\eta^2+s^2),\\
 \mathfrak b_R(\bar a)&\leq_*(64/3)Tv^2(\eta^2+s^2).
 \end{split}\label{eq:Omajorants}
\end{equation}
Here $\leq_*$ compares absolute coefficients as polynomials in $s,v^2$.

The entire signed term is bounded by taking $a_0=2\kappa$ in
\eqref{eq:oppositepurchase} and adding $2\kappa V^2\Qh^2$.
Indeed
\[
 2\kappa V^2\Qh^2+2\kappa\Rh^2
   -\kappa(V\Qh+\Rh)^2=\kappa(V\Qh-\Rh)^2 .
\]
The four fractions in this bound have majorants
\begin{equation}
 (12T+48Ts,\quad18v^2,\quad96Ts,\quad6T).
 \label{eq:signedmajorants}
\end{equation}
To see this directly, use the four bounds in \eqref{eq:kappabounds}
in \eqref{eq:oppositefractions}; for the $Q$ fraction use
$[4\phi\theta+2]\kappa/(q_cy)\leq18$.

\subsection*{The radial row}
The positive polynomial $\W$ and its permuted version $\W^\pi$ are given
in Appendix~\ref{app:parent}. Define
\[
 \mathcal R_A=\W+4U^2\Aop,\qquad
 \mathcal R_B=2U^2\W^\pi+
       [U\Zc(U^2-1)-(U^2-1)\Bh]^2 .
\]
They are sums of squares on the positive state chart.
Let
\begin{align}
 r_\Ps&=(1-\phi)\frac{4x^2y}{f_c}
       +\phi\left[\frac{8\theta xys}{f_c}
                     +\frac{4(x^2+xs+s^2)}{\Sigma}\right],\notag\\
 r_Q&=(1-\phi)\frac{4xys}{q_c}
        +\phi\frac{8\theta ys^2+4(1-\theta)xys}{q_c},\notag\\
 r_E&=\phi\frac{8s^2+2(s-x)^2}{\Sigma},\qquad
 r_R=(1-\phi)\frac{4xys}{r_c}.
 \label{eq:radialfractions}
\end{align}
Then the exact identity is
\begin{align}
 f_cr_\Ps\Ps+q_cr_Q\Qh^2+M_er_E\bc^2+r_cr_R\Rh^2-R_u^2
 =x^2y[(1-\phi)\W+\phi(1-\theta)\mathcal R_A
                                      +\phi\theta\mathcal R_B].
 \label{eq:radialpurchase}
\end{align}
Expansion using \eqref{eq:elementary} and
$Vw(U^2-1)=U\Zc(U^2-1)+(U^2-1)\Bh$ proves the identity.
The following coefficientwise bounds are sufficient:
\begin{equation}
 \begin{split}
 r_\Ps&\leq_*32B+(256/3)s+(512/3)(\eta+s)^2,\\
 r_Q&\leq_*64Bs+(256/3)s^2,\qquad
 r_E\leq_*(1280/3)(\eta+s)^2,\qquad
 r_R\leq_*(128/3)s .
 \end{split}\label{eq:radialmajorants}
\end{equation}
For example $x^2y/f_c\leq8B$, $xy/f_c\leq8/j$ and $y/q_c\leq8/j$.
Apply $\phi/j\leq4/3$ to the terms containing $s$ or $s^2$,
$\phi/\Sigma\leq128/3$ to the remaining terms, and
$(1-\phi)/(1-p)\leq4/3$ to the $r_R$ term.
For $(s-x)^2$, the coefficient majorant is $(s+x)^2$; their difference
$4xs$ is itself a positive polynomial square.

The bound \eqref{eq:Jpurchase} for $16J_\alpha^2$ includes
$32(1-\phi)(\alpha/B)^2R_u^2$.
Thus the effective radial multiplier is the source-only number
\begin{equation}
 L_R=16+32(1-\phi)(\alpha/B)^2
          \leq16+32\eta^2<17 .
 \label{eq:radialmultiplier}
\end{equation}
Its remaining $f_c\Ps$ and $r_c\Rh^2$ fractions are
\begin{equation}
 j_\Ps=\frac{16\phi j^2\alpha^2y(1+s/x)}{f_c},\qquad
 j_R=\frac{32(1-\phi)\alpha^2q_zs}{r_c}.
 \label{eq:Jfractions}
\end{equation}
Their majorants are $16[8B\eta^2+(32/3)\eta^2s]$ and
$16(64/3)\eta^2s$, respectively.

\subsection*{The two remaining helpers}
Put $\gamma=p^2/\Sigma\leq32$. From \eqref{eq:elementary},
\begin{align}
 \gamma M_ev^4\bc^2+\gamma f_c(yv^2+v^4)\Ps-H_1^2
       &=\gamma f_cyv^2\E,\notag\\
 \gamma TM_esv^2\bc^2+\gamma Tf_cs(y+v^2)\Ps-H_2^2
       &=\gamma Tf_cys\E .
 \label{eq:Hhelpers}
\end{align}
Both are polynomial identities, including on the testing axes.
The total fractions for $8H_1^2+8H_2^2$ are therefore
\[
 h_\Ps=8\gamma[yv^2+v^4+Ts(y+v^2)],\qquad
 h_E=8\gamma[v^4+Tsv^2].
\]

\subsection*{All scalar coefficients at once}
There are two equivalent ways to retain the signed square.
The complete-parent decomposition \eqref{eq:coreidentity} has already
used its exact remainder $\kappa(V\Qh-\Rh)^2$, leaving
$2\kappa V^2\Qh^2+2\kappa\Rh^2$ to be bounded.
Define its fractions by \eqref{eq:oppositefractions} with $a_0=2\kappa$,
plus $2\kappa V^2/q_c$ in the $Q$ component. Then the exact total losses are
\begin{equation}
 \begin{pmatrix}b_\Ps\\b_Q\\b_E\\b_R\end{pmatrix}
 =
 L_R\begin{pmatrix}r_\Ps\\r_Q\\r_E\\r_R\end{pmatrix}
 +\begin{pmatrix}j_\Ps\\0\\0\\j_R\end{pmatrix}
 +16\begin{pmatrix}\mathfrak b_\Ps(\bar a)\\
                    \mathfrak b_Q(\bar a)\\
                    \mathfrak b_E(\bar a)\\
                    \mathfrak b_R(\bar a)\end{pmatrix}
 +\begin{pmatrix}\mathfrak b_\Ps(2\kappa)\\
 \mathfrak b_Q(2\kappa)+2\kappa V^2/q_c\\
 \mathfrak b_E(2\kappa)\\\mathfrak b_R(2\kappa)\end{pmatrix}
 +\begin{pmatrix}h_\Ps+g_c/f_c\\0\\h_E\\0\end{pmatrix}.
 \label{eq:totallosses}
\end{equation}
Every replacement in this equality has the exact positive
remainder displayed above; the last $g_c/f_c$ bounds $g_c\Ps$ directly.
The radial term from \eqref{eq:Jpurchase} is included in $L_R$.

For explicit rational upper bounds, set
\[
 S=u_0^2=2^{-24},\qquad V_0=v_0^2=2^{-50},\qquad
 e=2\delta,\qquad D_0=e^2+S^2.
\]
Apply the nonnegative coefficient majorants at $s=S,v^2=V_0$.
The four resulting sums are bounded by
\begin{align}
 A_\Ps={}&17[32\delta+256S/3+512(e+S)^2/3]
 +16[8\delta e^2+32e^2S/3]\notag\\
 &+(2048/3)\delta V_0(1+2S)D_0
 +12\delta+48\delta S \notag\\
 &+256(\delta^2V_0+V_0^2+\delta^3S+\delta SV_0)
 +2(\delta^2+V_0)<9/16,\notag\\
 A_Q={}&17(64\delta S+256S^2/3)+512V_0^2D_0+18V_0<2^{-22},\notag\\
 A_E={}&(21760/3)(e+S)^2+(8192/3)\delta SV_0D_0
       +96\delta S+256(V_0^2+\delta SV_0)<1/16,\notag\\
 A_R={}&(2176/3)S+(1024/3)e^2S
                +(1024/3)\delta V_0D_0+6\delta<1/64 .
 \label{eq:rationaltotals}
\end{align}
These are strict rational comparisons at the displayed powers of two.
The last term in $A_\Ps$ is a conservative bound on $g_c/f_c\leq2y$.
Thus \eqref{eq:rationaltotals} proves \eqref{eq:totalbounds} as coefficient
norm bounds, not only as pointwise inequalities. In particular it proves
positivity of the fixed scalar reserves in the main text.


\section{The complete parent and its positive remainders}\label{app:parent}
For a direct expansion of \eqref{eq:slack}, put
\[
 q=v\zeta-Tuw=-T\sqrt x\,\Qh,\qquad
 m_p=\frac{\mu p}{1-p},
\]
and use the two independent normals $N,M$ in \eqref{eq:normal}.
Let
\begin{align}
 F_s={}&q_c\Qh^2+nj^2y\Bh^2+r_c\Rh^2+d_cV^2\Qh^2
 +(np^2y^2+\half p^2y^2-g_c)V^2\Bh^2\notag\\
 &+f_c\Ps+g_cV^2\Ps+k_c\Zc^2(V^2-1)^2,\notag\\
 L_N={}&nj\bc+npv^2\bc-j\alpha\bc+\half pv^2\bc
                         +K_RR_u+(K_c/j)uvq,\notag\\
 L_M={}&C_qq+\mu uv\bc+v(\alpha-BTCHu^2)R,\notag\\
 H_N={}&h_0(1+u^2)+\half v^2+\alpha u^2+n(1+v^2),\notag\\
 H_M={}&(1+u^2)[C(1+v^2)+D+\half]+A(1+v^2)+(\alpha+\half)v^2 .
 \label{eq:parentcoeff}
\end{align}
The full parent is the polynomial identity
\begin{equation}
 D_vS_c=F_s+2L_NN+2L_MM+H_NN^2+H_MM^2+
              2(\alpha+\half)uvNM+m_p(1+u^2)N^2 .
 \label{eq:fullparent}
\end{equation}
It follows by inserting \eqref{eq:rotation} into \eqref{eq:slack} and
using \eqref{eq:polys}--\eqref{eq:source}.
Equivalently these are the homogeneous RFN8/RFN12 identities of the
preserved companion source \texttt{latex/certificates/c2\_short.tex}.
The four scalar cancellations needed in this expansion are, with
$K_d=CH-(D+\half)(1-H)$,
\begin{align*}
 0&=T(K_d+AH+\half)-A,\\
 0&=jK_d+\mu+xTCH,\\
 0&=jK_d+xTCH+T[CH+AH+(\alpha+\half)H]
                   -A-(\alpha+\half)(1-p)+\alpha/B,\\
 0&=np-\mu+p/2+[-\mu-(\alpha+\half)(1-p)]x+j(\alpha-n)+yK_c .
\end{align*}
They are exact rational identities in the printed source dictionary.

The retained positive term in \eqref{eq:coreidentity} is explicitly
\begin{align}
 \calP={}&m_p(1+u^2)N^2+(h_0+\alpha)(uN)^2
             +\half(vN)^2+n[v(N+p\bc)]^2\notag\\
 &+(C+D+\half)(uM)^2+(C+A+\alpha+\half)(vM)^2+C(uvM)^2\notag\\
 &+d_c(V\Qh)^2+\half p^2v^2\bc^2
            +g_c[V(w-UV\Zc)]^2+\kappa(V\Qh-\Rh)^2
            +g_c(\Qh^2+\Rh^2).
 \label{eq:retainedall}
\end{align}
In particular $d_c=H^2(A+C+\alpha+\half)y^2q_z+\kappa>0$.
Combining
\[
 h_0N^2+n(N+j\bc)^2=\widehat N^2+E_*^2,\quad
 g_cV^2(\Ps-\Bh^2)=g_c(\E-\Ps),\quad
 -\kappa(V\Qh+\Rh)^2
 =\kappa(V\Qh-\Rh)^2-2\kappa V^2\Qh^2-2\kappa\Rh^2
\]
with \eqref{eq:fullparent} gives precisely \eqref{eq:coreidentity}
and the complete cross terms \eqref{eq:cross}.
Thus the decomposition is an identity for all four state coordinates.

\subsection*{The two radial sums of squares}
The radial polynomials used in \eqref{eq:radialpurchase} are
\begin{align}
 \W&=4\Ps+4U^2\Qh^2+4U^2\Rh^2-V^2(U^2-1)^2w^2,\notag\\
 \W^\pi&=4\Ps+4U^2\Qh^2+4U^2\Bh^2-\Zc^2(U^2-1)^2 .
 \label{eq:radialpolys}
\end{align}
Here is their complete positive Gram representation on the chart
$U,V\geq0$, $w\Zc\geq0$.
Put
\[
 t_7=\begin{pmatrix}
 \Zc-U^2Vw\\Vw-U^2Vw\\Uw-U^2Vw\\U\Zc-U^2Vw\\
 UVw-U^2Vw\\UV\Zc-U^2Vw\\U^2w-U^2Vw
 \end{pmatrix},\qquad
 t_3=\begin{pmatrix}1-UV\\V-UV\\U-UV\end{pmatrix}
\]
and
\[
 t_7^\pi=\begin{pmatrix}
 Vw-U^2\Zc\\\Zc-U^2\Zc\\Uw-U^2\Zc\\UVw-U^2\Zc\\
 U\Zc-U^2\Zc\\UV\Zc-U^2\Zc\\U^2w-U^2\Zc
 \end{pmatrix},\qquad
 t_3^\pi=\begin{pmatrix}w-U\Zc\\\Zc-U\Zc\\Uw-U\Zc\end{pmatrix}.
\]
With the following two fixed rational matrices,
\[
 G_7=\begin{pmatrix}
4&-2&-3/2&0&-1/2&0&1/2\\
-2&3&0&-3/2&0&-1&2\\
-3/2&0&4&3/2&-2&-3&0\\
0&-3/2&3/2&4&-3/2&0&-3/2\\
-1/2&0&-2&-3/2&3&2&-1\\
0&-1&-3&0&2&8&-4\\
1/2&2&0&-3/2&-1&-4&4
\end{pmatrix},\qquad
 G_3=\begin{pmatrix}3&-2&-2\\-2&2&2\\-2&2&3\end{pmatrix},
\]
the identities are
\begin{align}
 \W&=t_7^TG_7t_7+4Vw\Zc(1-U)^2+Uw\Zc\,t_3^TG_3t_3,\notag\\
 \W^\pi&=(t_7^\pi)^TG_7t_7^\pi+4Vw\Zc(1-U)^2
                                      +UV(t_3^\pi)^TG_3t_3^\pi .
 \label{eq:W19}
\end{align}
Their leading principal minors are
$(4,8,101/4,1049/16,907/16,3279/16,743/16)$ and $(3,2,2)$.
Both matrices are therefore positive definite. These explicit polynomial
identities are the homogeneous RFN19 and permuted RFN19 identities.
All extra square roots have fixed testing arguments.

For $w\Zc\leq0$, one instead uses the following direct expansions:
\begin{align}
 \W={}&(4U^2+3V^2+4U^4+3U^4V^2+2U^2V^2)w^2
 +(4+4U^2+8U^2V^2)\Zc^2
             -16UV(1+U^2)w\Zc,\notag\\
 \W^\pi={}&(4U^2+4V^2+4U^4+4U^2V^2)w^2
 +(3+2U^2+3U^4+8U^2V^2)\Zc^2
             -16UV(1+U^2)w\Zc .
 \label{eq:negativeW}
\end{align}
Every displayed monomial is a nonnegative coefficient times a square,
using the fixed radical $\sqrt{-UVw\Zc}$ for the last terms.
Together with the five-factor expansion of $\Ps$ in the main text, these
give a direct sum-of-squares identity also on the negative state chart.

\subsection*{The state signs and testing axes}
Let $K_d^{\mathrm{test}}$ denote the $d=1$ core in the numerator
of \eqref{eq:slack}, put $h=1+u^2+v^2$,
$\Delta=h^2-u^2v^2>0$, and set
\[
 A_0=\half\begin{pmatrix}h&-uv\\-uv&h\end{pmatrix},\qquad
 B_0=\half\begin{pmatrix}-u^2v&-uv^2\\uv^2&u^2v\end{pmatrix}.
\]
The exact core completion is
\begin{align}
 K_d^{\mathrm{test}}/D_u
 ={}&\left[\binom tr+A_0^{-1}B_0\binom w\zeta\right]^TA_0
 \left[\binom tr+A_0^{-1}B_0\binom w\zeta\right]\notag\\
 &+\frac{D_uD_v}{2\Delta}
      [u^2v^2(w^2+\zeta^2)-2uvh\,w\zeta] .
 \label{eq:signcore}
\end{align}
Since $A_0\succ0$, a negative core requires $uvw\zeta>0$.
Each of the nine original source summands is invariant under the gauges
\[
 (u,v,W,X,Y,Z)\longmapsto
 (\varepsilon u,\eta v,W,\eta X,\varepsilon Y,\varepsilon\eta Z).
\]
Thus a negative core can always be reduced to $u,v>0,w\zeta>0$.
At $u=0,v=0,w=0$ or $\zeta=0$, the remainder in \eqref{eq:signcore}
is nonnegative. The positive-chart proof and direct nonnegativity on the
complementary state set establish the theorem on the closed cap.


\IfFileExists{Fixed_Testing_Functions.tex}{\section{A fixed testing family and its affine equations}\label{app:functions}
The preceding identities give an evaluated finite certificate cone.
This section specifies its testing functions, including both state signs,
and proves that the cone is closed. Its admission statement is precisely
Theorem~\ref{thm:cap}; no complementary testing region is included here.

First take $u,v\geq0$ and use the common state sign to make $w\geq0$.
There are two closed charts, $\sigma w\zeta\geq0$, with $\sigma=+1,-1$.
All functions below are divided by $\sqrt{D_v}$. Their squared arguments
are the following fixed polynomials:
\[
 q_0=1,\quad q_1=uv,\quad q_2=uw\zeta,\quad q_3=vw\zeta
       \quad(\sigma=+1),\qquad
 q_0=1,\quad q_1=uv,\quad q_4=-uvw\zeta
       \quad(\sigma=-1).
\]
These weights are nonnegative on their respective charts. All source
radicals in a row are coefficients; for example
$\sqrt{UV}=x^{-1/4}y^{-1/4}\sqrt{uv}$ and
$\sqrt{Uw\Zc}=x^{-1/4}q_z^{-1/4}\sqrt{uw\zeta}$.
Thus no source parameter occurs in a testing-function argument.

\subsection*{A finite support prescription}
For a row $f$ with just one testing radical $\sqrt{q_k}$, let
$\operatorname{supp}_k f$ be the monomials in $u,v,t,r,w,\zeta$ that
remain after removing that radical and expanding. In taking this
support, source-only coefficients are formal symbols: possible
specializations to zero never delete a monomial. This convention makes
the family independent of the source, including at $n=0$.

Here is a complete finite prescription for the rows whose supports
are used. The ten noncircuit rows of $f$ in the main text are included,
as are the three circuit factors on the positive chart or the five
factors
\[
 Uw,\quad Vw,\quad\Zc,\quad UV\Zc,\quad2\sqrt{-UVw\Zc}
\]
on the negative chart. The latter identity is
$\Ps=U^2w^2+V^2w^2+\Zc^2+U^2V^2\Zc^2-4UVw\Zc$.
The retained squares in Appendix~\ref{app:parent} have row supports
\begin{gather*}
 N,\ uN,\ vN,\ v(N+p\bc),\ uM,\ vM,\ uvM,\quad
 V\Qh,\ V\Bh,\ V(w-UV\Zc),\ V\Qh-\Rh,\ \Qh,\ \Rh .
\end{gather*}
Every coefficient of these squares is nonnegative on the source domain.

Let $\mathbf A,\mathbf B,\mathbf E$ be the displayed square-root
vectors of $\Aop,\Bop,\E$ in \eqref{eq:elementary}--\eqref{eq:opposite};
let $\mathbf B_0=(\Rh,U\Qh,U(w-UV\Zc))$.
Let $\mathbf W_\sigma,\mathbf W^\pi_\sigma$ be the vectors in
\eqref{eq:W19} and \eqref{eq:negativeW}; on the positive chart their
constant positive blocks are $G_7,1,G_3$. Include every row of
\begin{gather*}
 \mathbf W_\sigma,\quad U\mathbf A,\quad
 U\mathbf W^\pi_\sigma,\quad (U^2-1)(U\Zc-\Bh),\quad
 (\alpha/B)R_u-\alpha uR,\quad\mathbf B_0,\\
 v\mathbf A,\quad u^2v\mathbf A,\quad
 v\mathbf B,\quad u^2v\mathbf B,\quad
 \sqrt T\,v(\alpha+c_0u^2)R,\quad
 \mathbf A,\quad\mathbf B,\quad v\mathbf E,\quad u\mathbf E,
 \qquad c_0=BTCH.
\end{gather*}
This list contains every positive remainder of the radial, $J_\alpha$,
$O$, signed-mass and helper identities. In particular, $\bar a$ is
split into its $v^2$ and $u^4v^2$ terms before taking roots; no
$\sqrt{\alpha^2+c_0^2u^4}$ is introduced.

The eight event coordinates are represented by the following rows,
with the squared rows in each group summed:
\[
\begin{array}{c|l}
 \lambda&W+vX\\
 n&X,\ vX\\
 m&vw-t,\ u(vw-t)\\
 C&r,\ ur,\ vr,\ uvr\\
 A&Y,\ vY\\
 D&r+v\zeta,\ u(r+v\zeta)\\
 \alpha&vY-Z\\
 \beta&t,\ ut,\ vt,\ uvt
\end{array}
\]
Here $W,X,Y,Z$ are the testing polynomials in \eqref{eq:rotation}.
These rows allow arbitrary nonnegative event additions literally.

For the four residual pivots take the rows of $\Ps,\Qh^2,\bc^2,\Rh^2$,
respectively. Their nonconstant loss supports are
\[
\begin{array}{c|l}
 \Ps&(0,2),(0,4),(2,0),(2,2),(4,0),(4,2),(6,2)\\
 Q&(0,2),(0,4),(2,0),(4,0),(4,4)\\
 E&(0,4),(2,0),(2,2),(4,0),(6,2)\\
 R&(0,2),(2,0),(4,2).
\end{array}\tag{\(*\)}
\]
For each pivot row $f$ and each $(a,b)$ in its row of $(*)$, include
$u^{a/2}v^{b/2}f$ in the base supports. Also form a separate block with
rows $\sqrt{1-P_{ab}}\,f$, where
$P_{ab}=(2^{12}u)^a(2^{25}v)^b$.
These rows implement both coefficient signs in the residual identity
in the main text. Constant zero coefficients may be retained.

Finally the two signs of the three matrix gates have respective
polynomial supports
\[
 \begin{array}{c|c|l}
 (a,b)&\hbox{weight}&\hbox{monomial vector}\\
 (1,1)&1\pm P_{11}&(t,r,vw,u\zeta,v\zeta,uw)\\
 (0,2),(2,2)&1\pm P_{ab}&(r,v\zeta,uw).
 \end{array}
\]
All rows participating in each nonconstant core coefficient lie in
the indicated span; the circuit rows do not participate.

For each chart and each radical class, collect the union of its base
supports into one monomial vector. For each scalar gate collect the
supports belonging to that exponent and class into a separate vector.
The matrix gates are the six vectors just displayed. A block with
vector $m_b$ and weight $q_b$ represents
\[
        D_v^{-1}q_b\,m_b^TQ_bm_b,\qquad Q_b\succeq0.
\]
This support prescription is an explicit finite construction, using
only polynomial expansion and collection of distinct monomials.
Its evaluated vectors and rational weights are in
\texttt{fixed\_family.json}; the full row list is in
\texttt{root\_catalog.json}. They give
\[
\begin{array}{c|rrr|r}
 \hbox{chart}&\hbox{base}&\hbox{scalar gates}&\hbox{matrix gates}&\hbox{sum}\\
 +&63&56&24&143\\
 -&58&49&24&131.
\end{array}
\]
There are $49$ Gram blocks, $2269$ upper-triangular Gram variables,
and total positive-semidefinite order $143+131+9=283$, including the
nine scalar source-sign blocks. A constant positive Gram block such
as $G_7$ is transferred to its monomial vector by a congruence with
source-only coefficients. Sums of such contributions remain positive.
Thus grouping rows does not change any identity or require a mixed inverse.

\subsection*{The rational affine system and its scope}
Let $T_c=D_vS_c$ after \eqref{eq:rotation}. Its nine coefficient columns are
\begin{gather*}
 T_\lambda=(W+vX)^2,\quad T_n=D_vX^2,\quad
 T_m=D_u(vw-t)^2,\quad T_C=D_uD_vr^2,\quad T_A=D_vY^2,\\
 T_D=D_u(r+v\zeta)^2,\quad T_\alpha=(vY-Z)^2,\quad
 T_\beta=D_uD_vt^2,\qquad
 T_d=(T_D+T_\alpha+T_\beta-D_vZ^2)/2.
\end{gather*}
For each chart impose the polynomial coefficient equations
\[
       \sum_{b\ {
m in\ chart}}q_bm_b^TQ_bm_b
          =\sum_{i=1}^9 c_iT_i .
\]
The file \texttt{affine\_system.json} records these $554$ equations
($277$ per chart): $2707$ Gram coefficients and $132$ source coefficients.
Indices start at zero; off-diagonal entries are actual symmetric Gram
entries and hence contribute twice. Every record is rational.

These equations and positive-semidefinite constraints define a cone
$K_{283}$ in the nine source coefficients. The displayed identities
and coefficient reserves prove that it contains every native source
in \eqref{eq:domain}, including $n=0$, and the eight event rays.
Conversely, each of its sources has $S_c\geq0$ for all states on the
closed testing cap: all block weights are nonnegative there.
To cover signed angles, apply
\[
 (u,v,t,r,w,\zeta)\longmapsto
 (a u,b v,t,ab r,bw,a\zeta),\qquad a,b\in\{-1,1\}.
\]
All nine $T_i$ are invariant. Reducing the angles to nonnegative
values and then fixing the common state sign reaches one of the two
charts. The testing axes and $w\zeta=0$ require no division and are
included by the same formulas. This is a certificate cone for the
stated cap, rather than a claim of adequacy on the remaining tests.

\subsection*{A fixed positive grid proves closedness}
Every monomial in $m_b$ has individual degree at most
$(4,3,1,1,1,1)$ in $(u,v,t,r,w,\zeta)$.
On either chart choose the $320$ rational points in
\[
 \left\{\frac{i u_0}{6}:1\leq i\leq5\right\}
 \times\left\{\frac{k v_0}{5}:1\leq k\leq4\right\}
 \times\{1/3,2/3\}^{3}\times\sigma\{1/3,2/3\}.
\]
All weights $q_b$ are strictly positive at these points. Distinct
monomials with the displayed individual degree bounds are independent
on this Cartesian grid by successive univariate interpolation.
Consequently the rational matrix
\[
       H_b=\sum_{a\ \rm in\ grid}
              \frac{q_b(a)}{1+v(a)^2}\,m_b(a)m_b(a)^T
\]
is positive definite for every block.

Suppose $c^{(k)}\in K_{283}$ converges to $c$. Evaluate its identity
on these finitely many points and sum. Each summand from each Gram
block is nonnegative, so $\operatorname{tr}(H_bQ_b^{(k)})$ is bounded
above by the sum of the target values. These target sums converge
because they depend linearly on $c^{(k)}$. Since $H_b\succ0$, every
$Q_b^{(k)}$ is bounded. There are finitely many blocks, so a common
subsequence converges to positive-semidefinite matrices satisfying
the same coefficient equations at $c$. Hence $K_{283}$ is closed.
This argument also handles dependencies between different blocks;
it needs only independence within each explicitly chosen monomial vector.
}{
 \section{Fixed testing functions}\label{app:functions}
 The explicit testing-function lists and their coefficient prescription are
 supplied in the separate input \texttt{Fixed\_Testing\_Functions.tex}.}
\end{document}

